\section{Explicit principal-cover gluing}

\subsection{Statement}

Let $A$ be a commutative ring, $M$ an $A$-module, and suppose
\[
                  \sum_{i<r}c_i f_i=1.
\]
For each $i$, let a local section on $D(f_i)$ be represented by
\[
                  s_i=\frac{m_i}{f_i^{n_i}}\in M_{f_i}.
\]
On the overlap $D(f_if_j)$, the two restrictions are represented after denominator
homogenization by fractions in the single localization $M_{f_if_j}$.

\begin{theorem}[Primitive-recursive Zariski gluing]\label{thm:principal-gluing}
Assume that for every $i,j<r$ a finite localization witness proves that $s_i$ and $s_j$ agree in
$M_{f_if_j}$.  Then an explicit primitive-recursive formula produces $m\in M$ and witnesses that
$m/1=s_i$ in every $M_{f_i}$.  If two global sections have the same restrictions, an explicit
primitive-recursive witness proves them equal.  The theorem is valid over rings with zero divisors.
\end{theorem}

This is the sheaf equalizer theorem for the quasi-coherent module associated to $M$, presented in a
form directly usable by iPRA.

\subsection{Denominator homogenization}

Set
\[
                         N=\max_{i<r} n_i,
          \qquad a_i=f_i^{N-n_i}m_i.
\]
Then $s_i=a_i/f_i^N$.  An overlap witness can be written as an exponent $k_{ij}$ satisfying
\begin{equation}\label{eq:overlap-kill}
  (f_if_j)^{k_{ij}}
  \bigl(f_j^N a_i-f_i^N a_j\bigr)=0_M.
\end{equation}
The witness includes the module-equality certificate for this equation.  Let
\[
                         K=\max_{i,j<r}k_{ij},
             \qquad q=N+K+1,
             \qquad b_i=f_i^{q-N}a_i=f_i^{K+1}a_i.
\]
Multiplying \eqref{eq:overlap-kill} by the remaining powers needed to reach $K$ gives
\begin{equation}\label{eq:cross-common}
                         f_j^q b_i=f_i^q b_j
                         \qquad(i,j<r).
\end{equation}
Indeed the difference of the two sides equals
\[
 f_if_j(f_if_j)^K(f_j^Na_i-f_i^Na_j),
\]
which is zero by the synchronized overlap witness.  The harmless extra $+1$ in $q$ avoids a special
case when every denominator and every killing exponent is zero.

\subsection{Construction of the global section}

Apply Theorem~\ref{thm:power-cover} to the exponent $q$.  It returns $d_i$ satisfying
\[
                         \sum_{i<r}d_i f_i^q=1.
\]
Define
\begin{equation}\label{eq:global-glue}
                         m=\sum_{i<r}d_i b_i.
\end{equation}
For a fixed chart $j$, use \eqref{eq:cross-common}:
\begin{align*}
 f_j^q m
   &=\sum_i d_i f_j^q b_i\\
   &=\sum_i d_i f_i^q b_j\\
   &=b_j.
\end{align*}
Thus $m/1=b_j/f_j^q=a_j/f_j^N=m_j/f_j^{n_j}$ in $M_{f_j}$.  Every equality in this calculation is
accompanied by a finite witness obtained from module distributivity, the pairwise cross witnesses,
and the power-cover certificate.

This proves existence without search.  Once the overlap witnesses and the unit-ideal cover are in
hand, the global section is a fixed bounded formula.

\subsection{Uniqueness}

Suppose $m,m'\in M$ have equal restrictions in every $M_{f_i}$.  For each $i$, let $e_i$ witness
\[
                            f_i^{e_i}(m-m')=0.
\]
Set $E=1+\max_i e_i$.  Then $f_i^E(m-m')=0$ for every $i$.  Use the power-cover constructor to find
$u_i$ with $\sum u_i f_i^E=1$.  Therefore
\[
  m-m'=\sum_i u_i f_i^E(m-m')=0.
\]
This is a direct equality witness in $M$.  No cancellation, prime ideal argument, or excluded middle
is used.

\begin{corollary}[Finite principal-cover equalizer]\label{cor:equalizer}
The sequence
\[
  0\longrightarrow M
   \longrightarrow \prod_{i<r}M_{f_i}
   \rightrightarrows \prod_{i,j<r}M_{f_if_j}
\]
is exact in the following certificate sense:
\begin{enumerate}[label=(\alph*)]
\item the first map has an explicit joint-monicity witness transformer;
\item every tuple with supplied pairwise compatibility witnesses has an explicitly constructed
      preimage;
\item the restriction witnesses and uniqueness witness are primitive recursive in the input codes.
\end{enumerate}
\end{corollary}

Finite products are coded by lists, so the theorem is a statement entirely about natural numbers.
The use of the word ``exact'' does not import an abstract abelian-category library.

\subsection{A zero-divisor example}

Take $A=M=\Z/12\Z$, $f_0=3$, $f_1=4$, and cover coefficients $c_0=3$, $c_1=1$:
\[
                    3\cdot3+1\cdot4=13=1\pmod {12}.
\]
The local numerators $9$ on $D(3)$ and $8$ on $D(4)$ both represent the global class $5$:
\[
                  3(9-5)=0,\qquad 4(8-5)=0\pmod {12}.
\]
On the overlap, $(3\cdot4)^1=0$ kills every difference, so the compatibility witness has exponent
one.  The gluing algorithm takes $N=0$, $K=1$, $q=2$ if the displayed formula is used literally.
A slightly optimized implementation uses $q=1$ because the synchronized cross identity already
holds; either choice is sound.  The executable companion uses the optimized exponent and returns the
unique global class $5$.

This example demonstrates why localization equality must retain killing exponents.  Cross
multiplication alone would not be a cancellative equality test in a ring with zero divisors.

\subsection{Primitive-recursive complexity}

All loops are bounded by primitive-recursive terms in the input lengths and exponents:

\begin{itemize}
\item scanning the $n_i$ and $k_{ij}$ uses bounded maxima;
\item exponentiation uses primitive recursion;
\item the power-cover constructor enumerates finitely many weak compositions or performs a bounded
      iterative expansion;
\item the final sums are finite module folds.
\end{itemize}

A convenient quantifier-free invariant for the final fold over $i<t$ is
\[
 f_j^q m_t=\left(\sum_{i<t}d_i f_i^q\right)b_j,
\]
where $m_t=\sum_{i<t}d_i b_i$.  The invariant is an equality of coded module terms with a concrete
witness accumulator.  It contains no quantification over local sections or opens.

\subsection{Refinements and iterated covers}

If a chart $D(f_i)$ is further covered by principal opens $D(f_i g_{i\alpha})$ with a local
unit-ideal certificate in $A_{f_i}$, clear the finitely many denominators in that certificate and
obtain a power cover in $A$.  Gluing first within each chart and then across the outer cover agrees
with gluing on the flattened refinement.  The proof follows uniqueness: both constructions have the
same restrictions on every refined principal piece, so the global outputs are equal.

Thus finite cover refinement is coherent without choosing canonical normal forms.  This is the
basic associativity needed to build finite affine atlases and sheaf data.
